\documentclass[11pt,a4paper]{article}

\usepackage[utf8]{inputenc}
\usepackage[T1]{fontenc}
\usepackage[british]{babel}
\usepackage[margin=2cm]{geometry}
\usepackage{amsmath}
\usepackage{booktabs}
\usepackage{graphicx}
\usepackage{array}
\usepackage{url}
\usepackage[hidelinks]{hyperref}

\newcommand{\eps}{\varepsilon}
\newcommand{\inst}[1]{\texttt{#1}}

\title{Dependency context, declaration slices and inert padding in repository-level patch generation:\\ exploratory ablations following a preregistered pilot}
\author{Jude Castiel Johansen\\ Shockware Limited (independent researcher)}
\date{}

\begin{document}
\maketitle

\begin{abstract}
This report describes three phases of exploratory ablations that follow a preregistered pilot \cite{pilot}. All three phases are exploratory. They ask three questions. First, does adding dependency documentation to a base context (arm A1 against arm W) change the rate at which sampled patches pass? Second, are reduced contexts that keep only predicted dependency groups non-inferior to the full dependency context? Third, can declaration-level slices and inert padding of matched length separate an effect of content from an effect of length? The evidence base is small. Phase 1 uses 23 instances from one repository. Phase 2 uses 10 instances from two repositories. Phase 3 uses eleven instances from two repositories (8 from \texttt{cli/cli}, 3 from \texttt{caddyserver/caddy}).

Phase 1 found no leave-one-out necessity signal beyond a simulated null. In Phase 2, most relative contrasts are near-vacuous because many instances sit at or near zero pass rate. In Phase 3, on Haiku 5.5, inert padding matched to A1's prompt tokens lowered pass rates relative to both W and A1, and both intervals exclude zero. A declaration slice of W (Wd) was non-inferior to W at $\eps = 0.05$ on the post-screen samples. It was not non-inferior on all samples, and the choice of sample set was not preregistered. On Sonnet 5.5, Wd was worse than W (the interval excludes zero), and 6/11 instances were at ceiling.

None of these results distinguishes an information effect from a reader effect with confidence.
\end{abstract}

\section{Scope and status}

This is an exploratory report. Phases 1 to 3 were run after the pilot \cite{pilot}, and their designs were amended during execution. The amendments and integrity events are listed in Section~\ref{sec:deviations}.

No result below is confirmatory. In particular, hypothesis H5, which the preregistration designates as confirmatory on the Phase 3 instances, was not tested confirmatorily (Section~\ref{sec:h5}).

The prior records in this series are \cite{prior} and \cite{pilot}. This report makes no claim of novelty beyond them.

The numbers reported here are taken from the verified results file generated by \texttt{analysis/phase23.py}. That file draws on the verdict records in \texttt{harness/results-abl}, \texttt{harness/results-p2}, \texttt{harness/results-p2screen} and \texttt{harness/results-p3}, and on the arm metadata in \texttt{harness/abl-*/} and \texttt{harness/ablr-*/contexts.json}.

\section{Methods}

\subsection{Instances}

Instances are mined commits with fail-to-pass tests. Full instance identifiers take the form \inst{cli\_\_cli-\textit{hash}} and \inst{caddyserver\_\_caddy-\textit{hash}}. In the text they are abbreviated as ``cli \textit{hash}'' and ``caddy \textit{hash}''.

The valid mined instances on disk, by repository prefix, are:
\begin{itemize}
\item \inst{caddyserver\_\_caddy}: 14;
\item \inst{cli\_\_cli}: 30;
\item \inst{gohugoio\_\_hugo}: 1;
\item \inst{golangci\_\_golangci-lint}: 1;
\item \inst{restic\_\_restic}: 1.
\end{itemize}

The \inst{cli\_\_cli} count includes the 24 pilot instances. One of these is invalid because its task text is empty, which leaves 23 valid pilot instances. Mining used commits since 2026-06-01.

\subsection{Arms}

\begin{itemize}
\item \textbf{W} is the base context.
\item \textbf{A1} adds the documentation of the instance's dependency groups to W.
\item \textbf{A1g} and \textbf{A1t} keep only the dependency groups selected by a gold-referenced predictor and a task-referenced predictor, respectively. When a predictor keeps no group, the arm equals W.
\item \textbf{A0} is used only for Opus 5.5 in Phase 2. It is defined in the preregistration.
\end{itemize}

Phase 3 adds four arms:
\begin{itemize}
\item \textbf{Wd} contains the declarations of W that are reachable, through the type-checked reference graph, from identifiers in the fail-to-pass tests and from task-text code spans and selectors. Imports and package clauses are kept, and omitted declarations are marked.
\item \textbf{A1r} is the full W plus only the dependency declarations that W references, rendered without bodies.
\item \textbf{Wdr} is the slice plus the dependency declarations that the slice references.
\item \textbf{Wpad} is W preceded by inert same-module padding drawn from the pilot pool. The padding is matched to A1's prompt tokens.
\end{itemize}

\subsection{Models and sample sizes}

The models are:
\begin{itemize}
\item \texttt{claude-haiku-5-5} (Haiku 5.5);
\item \texttt{claude-haiku-4-5-20251001};
\item \texttt{claude-sonnet-5-5} (Sonnet 5.5);
\item \texttt{claude-opus-5-5} (Opus 5.5).
\end{itemize}

Samples per arm and instance ($k$) are given with each phase.

\subsection{Analysis}

The outcome is the per-instance pass rate of an arm. Paired contrasts use an instance-level percentile bootstrap with 2,000 resamples and seed 2. The seed differs from the one in the original analysis plan. The change is part of the analysis amendment of 2026-10-12.

For each contrast we report the following:
\begin{itemize}
\item the mean difference;
\item the two-sided 95\% interval;
\item the one-sided 95\% lower bound (LB);
\item sign counts (+/$-$/0);
\item the number of instances $n$.
\end{itemize}

The decision bound is the one-sided 95\% LB, as set by the amendment of 2026-10-12. The lower end of the two-sided interval corresponds to a one-sided 2.5\% test, and it is not used as the decision bound. A contrast is called non-inferior at $\eps = 0.05$ only when the LB exceeds $-0.05$.

An interval that contains zero is reported as an unresolved difference. It is not reported as equivalence.

A relative comparison on an instance is informative only if the best arm on that instance reaches a pass rate of at least $0.1$ (the non-degeneracy condition). Instances whose best arm is below $0.1$ are counted as degenerate.

\subsection{Information effects and reader effects}

Three kinds of effect are distinguished here. The distinction is used only to organise the contrasts. None of the three has been established.

\begin{itemize}
\item An \emph{information effect} means that the added context supplies something the passing patch needs.
\item The first reader effect, a \emph{load effect}, means that the length of the prompt changes pass rates whatever its content. The contrasts Wpad $-$ W and Wpad $-$ A1 are designed to expose it. Wpad $-$ A1 is the matched-length comparison (median token ratio 1.00).
\item The second reader effect, a \emph{presentation effect}, means that the presence or form of dependency material changes pass rates without supplying anything the patch uses. The diagnostic of 2026-10-13 (Section~\ref{sec:diagnostic}) bears on this possibility. That diagnostic is not a script output and is not a test of it.
\end{itemize}

\section{Phase 1 (2026-10-12): Haiku 5.5 on 23 cli instances}

Phase 1 ran Haiku 5.5 at $k=20$ on the 23 valid pilot instances. It used the task-text rule of the original H5 paragraph. A build-environment fault invalidated the verdicts for 4,456 samples. Of these, 4,409 were re-verified from stored responses with a patch-hash equality check, and 44 were re-sampled.

The faulty aggregate was seen before the repair. For this reason Phase 1 is exploratory in every respect.

Two runs of A1 can be compared directly. A1 in Phase 1 ($k=20$) minus A1 in the pilot ($k=10$) is $-0.041$ $[-0.074, -0.013]$ over 23 instances.

Table~\ref{tab:p1} gives the contrasts. The counts that limit them are as follows:
\begin{itemize}
\item A1g equals W on 14/23 instances, and A1t equals W on 11/23.
\item W and A1 are both at ceiling (1.0) on 8 instances and both at floor (0.0) on 5.
\item W is inexpressible (gold identifiers missing from W) on 10 instances.
\end{itemize}
Only instances outside ceiling and floor can separate the arms. Per-instance absolute pass rates for Phase 1 are not reproduced in the results file used for this report.

\begin{table}[h]
\centering
\caption{Phase 1 paired contrasts (Haiku 5.5, $k=20$, $n=23$). ``Without most influential'' drops the single instance named. The token columns give the median prompt tokens of the first arm and of A1.}
\label{tab:p1}
\resizebox{\textwidth}{!}{%
\begin{tabular}{lcccccccc}
\toprule
contrast & mean & 95\% CI & LB & non-inferior & signs & without most influential & tokens & tokens A1\\
\midrule
A1g $-$ A1 & $+0.017$ & $[-0.002, +0.041]$ & $+0.000$ & yes & 5/2/16 & $+0.011$ $[-0.005, +0.032]$ (drop 2bcdfbb1ac02) & 13,513 & 39,586\\
A1t $-$ A1 & $+0.007$ & $[-0.020, +0.030]$ & $-0.015$ & yes & 5/3/15 & $+0.000$ $[-0.023, +0.023]$ (drop 2bcdfbb1ac02) & 15,335 & 39,586\\
W $-$ A1 & $-0.028$ & $[-0.096, +0.017]$ & $-0.083$ & no & 4/4/15 & $+0.000$ $[-0.025, +0.023]$ (drop babd691446b9) & 13,476 & 39,586\\
\bottomrule
\end{tabular}}
\end{table}

A1g $-$ A1 and A1t $-$ A1 are non-inferior at $\eps = 0.05$. Both intervals contain zero, so the differences are unresolved. They are not shown to be equivalent, and these results do not show that the arms make no difference.

W $-$ A1 is not non-inferior (LB $-0.083$). After dropping babd691446b9 the estimate is $+0.000$ $[-0.025, +0.023]$. That exclusion of a single instance is post hoc, and it is reported here only as a sensitivity check. It is not a finding.

\paragraph{Leave-one-dependency-out.}
On 17/23 instances A1 has a pass rate above 0, and these instances carry 135 dependency groups. A group was labelled ``necessary'' in 33/135 cases (0.244 $[0.178, 0.319]$, bootstrap over groups). This interval ignores instance clustering. Removing a group raised the pass rate by more than $\eps$ in another 33/135 cells.

Under a null in which all cells of an instance share one rate ($k=20$, 2,000 simulations), the expected counts are 30.8 $[13, 49]$ for ``necessary'' and 30.4 $[12, 53]$ for ``helps''. Both observed counts lie inside these ranges. They are therefore not evidence of necessity.

The preregistration's §8 declares precision and recall against this label uninterpretable. Because the label does not exceed its null, these quantities are not reported as measures of predictor quality.

\section{Phase 2 (2026-10-12/13): four models on 10 instances}

\paragraph{Screen.}
W was screened at $k=5$ on 35 instances (22 cli, 10 caddy, 1 hugo, 1 restic, 1 golangci-lint). On 16 of these W was inexpressible. Ten instances were selected (6 cli, 4 caddy), using the criteria ``W inexpressible'' and ``W $< 0.8$''.

The task-text rule is that of the 2026-10-12 amendment (selectors, code-span identifiers, import paths).

\paragraph{Arms and sample sizes.}
The dense models ran W, A1, A1g and A1t. Opus 5.5 ran A0, W and A1t. The values of $k$ were:
\begin{itemize}
\item \texttt{claude-haiku-5-5}: 100;
\item \texttt{claude-haiku-4-5-20251001}: 20;
\item \texttt{claude-sonnet-5-5}: 10;
\item \texttt{claude-opus-5-5}: 5.
\end{itemize}

\paragraph{Prompt tokens.}
Median prompt tokens on the 10 instances were W 21,946, A1 84,817, A1g 49,822 and A1t 49,822.

\paragraph{Missing cells.}
For \texttt{claude-haiku-4-5-20251001}, no cells are recorded for caddy 7910dd94eea2, and caddy 6727b2f3a92b has no A1 cell. Contrasts involving A1 therefore have $n=8$ for this model, and A1t $-$ W has $n=9$. For Opus 5.5, cli 2bcdfbb1ac02 has no A1t cell, so A1t $-$ W has $n=9$.

\begin{table}[h]
\centering
\caption{Phase 2 per-instance pass rates. An arm identical to another arm is shown with the rate of that arm. ``deps'' is the number of dependency groups. ``A1g/A1t kept'' is the number of groups each predictor keeps. A dash means no cell was recorded.}
\label{tab:p2inst}
\scriptsize
\begin{tabular}{l c c p{3.2cm} p{3.0cm} p{3.0cm} p{3.0cm}}
\toprule
instance & deps & A1g/A1t kept & \texttt{claude-haiku-5-5} & \texttt{claude-haiku-4-5-20251001} & \texttt{claude-sonnet-5-5} & \texttt{claude-opus-5-5}\\
\midrule
caddy 4f45ded75556 & 8 & 1/0 & W 0.03, A1 0.03, A1g 0.01, A1t 0.03 & all arms 0.00 & all arms 0.00 & A0 0.00, W 0.00, A1t 0.00\\
caddy 6727b2f3a92b & 13 & 0/0 & W 0.44, A1 0.22, A1g 0.44, A1t 0.44 & W 0.00, A1g 0.00, A1t 0.00 & W 0.20, A1 0.10, A1g 0.20, A1t 0.20 & A0 0.00, W 1.00, A1g 1.00, A1t 1.00\\
caddy 7910dd94eea2 & 12 & 1/0 & W 0.01, A1 0.07, A1g 0.02, A1t 0.01 & -- & all arms 0.00 & A0 0.60, W 0.00, A1t 0.00\\
caddy c8591a1e46e4 & 5 & 1/2 & W 0.57, A1 0.46, A1g 0.49, A1t 0.45 & W 0.40, A1 0.15, A1g 0.15, A1t 0.05 & W 0.20, A1 0.60, A1g 0.00, A1t 0.00 & A0 0.00, W 0.00, A1t 0.00\\
cli 015e6cf6ac39 & 0 & 0/0 & all arms 0.00 & all arms 0.00 & all arms 0.00 & all arms 0.00\\
cli 20147bf81fed & 50 & 3/0 & W 0.00, A1 0.01, A1g 0.00, A1t 0.00 & all arms 0.00 & all arms 0.00 & A0 0.00, W 1.00, A1t 1.00\\
cli 2bcdfbb1ac02 & 4 & 2/2 & W 0.00, A1 0.27, A1g 0.14, A1t 0.14 & all arms 0.00 & W 0.00, A1 0.10, A1g 0.10, A1t 0.10 & A0 0.80, W 0.80\\
cli 9b6585be5218 & 8 & 1/0 & W 0.01, A1 0.00, A1g 0.00, A1t 0.01 & all arms 0.00 & W 0.10, A1 0.00, A1g 0.00, A1t 0.10 & A0 0.00, W 0.00, A1t 0.00\\
cli 9ee0afb43414 & 10 & 0/0 & all arms 0.00 & all arms 0.00 & all arms 0.00 & all arms 0.00\\
cli ff28e72b59be & 62 & 0/0 & all arms 0.00 & W 0.00, A1 0.05, A1g 0.00, A1t 0.00 & W 1.00, A1 0.80, A1g 1.00, A1t 1.00 & A0 1.00, W 1.00, A1g 1.00, A1t 1.00\\
\bottomrule
\end{tabular}
\end{table}

\begin{table}[h]
\centering
\caption{Phase 2 paired contrasts. ``Degenerate'' is the number of instances whose best arm is below 0.1.}
\label{tab:p2}
\small
\begin{tabular}{llccccc}
\toprule
model & contrast & mean [95\% CI] & LB & non-inferior & signs & $n$\\
\midrule
\texttt{claude-haiku-5-5} & A1t $-$ A1 & $+0.002$ $[-0.043, +0.059]$ & $-0.038$ & yes & 2/4/4 & 10\\
(degenerate 7/10) & A1g $-$ A1 & $+0.004$ $[-0.042, +0.061]$ & $-0.035$ & yes & 2/4/4 & 10\\
 & W $-$ A1 & $+0.000$ $[-0.076, +0.071]$ & $-0.060$ & no & 3/3/4 & 10\\
 & A1t $-$ W & $+0.002$ $[-0.036, +0.042]$ & $-0.024$ & yes & 1/1/8 & 10\\
\midrule
\texttt{claude-haiku-4-5-20251001} & A1t $-$ A1 & $-0.019$ $[-0.044, +0.000]$ & $-0.037$ & yes & 0/2/6 & 8\\
(degenerate 9/10) & A1g $-$ A1 & $-0.006$ $[-0.019, +0.000]$ & $-0.019$ & yes & 0/1/7 & 8\\
 & W $-$ A1 & $+0.025$ $[-0.019, +0.094]$ & $-0.013$ & yes & 1/1/6 & 8\\
 & A1t $-$ W & $-0.039$ $[-0.117, +0.000]$ & $-0.117$ & no & 0/1/8 & 9\\
\midrule
\texttt{claude-sonnet-5-5} & A1t $-$ A1 & $-0.020$ $[-0.170, +0.080]$ & $-0.150$ & no & 3/1/6 & 10\\
(degenerate 5/10) & A1g $-$ A1 & $-0.030$ $[-0.180, +0.060]$ & $-0.150$ & no & 2/1/7 & 10\\
 & W $-$ A1 & $-0.010$ $[-0.120, +0.070]$ & $-0.100$ & no & 3/2/5 & 10\\
 & A1t $-$ W & $-0.010$ $[-0.060, +0.030]$ & $-0.050$ & no & 1/1/8 & 10\\
\midrule
\texttt{claude-opus-5-5} & A1t $-$ W & $+0.000$ $[+0.000, +0.000]$ & $+0.000$ & degenerate & 0/0/9 & 9\\
(degenerate 5/10) & A0 $-$ W & $-0.140$ $[-0.440, +0.120]$ & $-0.400$ & no & 1/2/7 & 10\\
\bottomrule
\end{tabular}
\end{table}

\paragraph{Interpretation.}
The Phase 2 relative contrasts are near-vacuous because of floor effects. The numbers of degenerate instances are 7/10 for Haiku 5.5, 9/10 for \texttt{claude-haiku-4-5-20251001}, 5/10 for Sonnet 5.5 and 5/10 for Opus 5.5.

For Haiku 5.5, A1t $-$ A1 is non-inferior at $\eps = 0.05$ (LB $-0.038$). This rests on the 3 instances that meet the non-degeneracy condition. Their absolute rates are as follows:
\begin{itemize}
\item caddy 6727b2f3a92b: W 0.44, A1 0.22, A1t 0.44;
\item caddy c8591a1e46e4: W 0.57, A1 0.46, A1t 0.45;
\item cli 2bcdfbb1ac02: W 0.00, A1 0.27, A1t 0.14.
\end{itemize}

The non-inferior verdicts for \texttt{claude-haiku-4-5-20251001} rest on a single non-degenerate instance (caddy c8591a1e46e4, W 0.40). They carry almost no information.

For Sonnet 5.5, A1t $-$ W has LB $-0.050$. This does not exceed $-0.05$, so the contrast is not non-inferior.

For Opus 5.5, A1t $-$ W has the interval $[+0.000, +0.000]$ with signs 0/0/9. This interval is degenerate, and it is not evidence of sufficiency.

\section{Leave-one-dependency-out at $k=100$ (Haiku 5.5)}

Stage 6 of Phase 2 was restricted on 2026-10-13 to the three instances with A1 $\geq 0.1$. The difference reported is pass(A1 minus group) $-$ pass(A1). The null reference draws all cells of an instance from one pooled rate at $k=100$ (2,000 simulations) and counts removals with $|\text{difference}| > 0.05$.

\begin{table}[h]
\centering
\caption{Leave-one-dependency-out at $k=100$. Observed counts of removals with $|\text{difference}| > 0.05$, with the null reference.}
\label{tab:loo}
\small
\begin{tabular}{lccccc}
\toprule
instance & A1 & W & tokens A1 / W & observed & null reference\\
\midrule
caddy 6727b2f3a92b & 0.22 & 0.44 & 277,505 / 158,565 & 9 of 13 & 5.5 $[1, 12]$\\
caddy c8591a1e46e4 & 0.46 & 0.57 & 84,817 / 12,137 & 2 of 5 & 2.3 $[0, 5]$\\
cli 2bcdfbb1ac02 & 0.27 & 0.00 & 16,293 / 2,367 & 2 of 4 & 1.5 $[0, 4]$\\
\bottomrule
\end{tabular}
\end{table}

Every observed count lies inside its null range. No necessity claim is supported at the count level.

Descriptively, on cli 2bcdfbb1ac02, removing the \texttt{queries} group moved the pass rate from 0.27 to 0.00 (difference $-0.27$, $n=100$). Removing \texttt{iostreams} gave 0.18. In the Phase 2 cell for this instance, the arm that keeps \texttt{queries} and \texttt{iostreams} (A1g = A1t) passed at 0.14, and W passed at 0.00.

On caddy 6727b2f3a92b, most removals raised the pass rate. For example, removing \texttt{caddytls} gave 0.44 (difference $+0.22$). This is consistent with W (0.44) exceeding A1 (0.22) on that instance, but the count of 9 of 13 is within its null range.

\section{Diagnostic of 2026-10-13 (not a script output)}
\label{sec:diagnostic}

A harness session on 2026-10-13 examined the three cells in which context changed the outcome:
\begin{itemize}
\item caddy c8591a1e46e4 on Sonnet 5.5: A1 6/10, W 2/10, A1g 0/10, A1t 0/10;
\item cli 2bcdfbb1ac02 on Haiku 5.5: A1 27/100, W 0/100;
\item caddy 6727b2f3a92b on Haiku 5.5: W 44/100, A1 22/100, with W at 158,565 tokens and A1 at 277,505 tokens.
\end{itemize}

The session recorded that the passing patches introduced no code identifier that only the dependency documentation supplied. The only tokens in passing patches that came solely from the dependencies were English words from comments.

This observation is not a script output and has not been reproduced. Attributing the differences to how the model reads the context, rather than to information supplied, is an untested mechanism. It is recorded here as an unverified interpretation.

\section{Phase 3 (2026-10-13/14): declaration-level arms and matched padding}

\paragraph{Screen.}
W and A1 were screened at $k=20$ on Haiku 5.5 over 35 instances. Twelve instances were selected because either rate lay in $[0.1, 0.9]$. These included gohugoio/hugo c9f6e4d4abc3, which was dropped on 2026-10-14 because its single instance produced a 182 GB build cache. Eleven instances were run: 8 cli and 3 caddy.

\paragraph{Arms and sample sizes.}
Haiku 5.5 ran at $k=50$ on W, A1, A1t, Wd, A1r, Wdr and Wpad. Its W and A1 arms continue the $k=20$ screen samples. Sonnet 5.5 ran at $k=10$ on W, A1, Wd, A1r, Wdr and Wpad.

Every instance meets the non-degeneracy condition: Haiku 11/11 and Sonnet 11/11. On Sonnet, W and A1 are both at 1.00 on 6/11 instances.

\paragraph{Prompt tokens.}
Median prompt tokens were W 22,957, Wd 16,439, A1 68,609, A1r 31,357, Wdr 16,776 and Wpad 68,718. Wd is byte-identical to W on cli 2bcdfbb1ac02, cli 326faaac8b5c and cli 8f7afbdbdbd6, because every declaration on those instances is reachable.

\begin{table}[h]
\centering
\caption{Phase 3 prompt tokens and declaration counts. ``slice decls'' is the number of declarations in the slice against the number in W. ``dep decls'' is the number of dependency declarations referenced by the slice against the number referenced by W.}
\label{tab:p3tok}
\small
\begin{tabular}{lrrrrrrcc}
\toprule
instance & W & Wd & A1 & A1r & Wdr & Wpad & slice decls & dep decls\\
\midrule
caddy 30d3bd2a8959 & 171,794 & 82,533 & 237,893 & 198,792 & 89,753 & 183,345 & 189/606 & 42/126\\
caddy 6727b2f3a92b & 158,565 & 73,923 & 277,505 & 203,263 & 93,699 & 229,911 & 181/459 & 68/157\\
caddy c8591a1e46e4 & 12,137 & 10,047 & 84,817 & 17,460 & 13,474 & 84,798 & 23/37 & 23/27\\
cli 2bcdfbb1ac02 & 2,367 & 2,367 & 16,293 & 4,443 & 4,443 & 16,341 & 7/7 & 14/14\\
cli 326faaac8b5c & 21,011 & 21,011 & 67,579 & 31,357 & 31,357 & 68,049 & 44/44 & 79/79\\
cli 3f6a16a9f8c7 & 13,485 & 4,045 & 24,999 & 17,612 & 5,744 & 24,960 & 6/17 & 10/43\\
cli 8f7afbdbdbd6 & 16,439 & 16,439 & 56,667 & 21,820 & 21,820 & 56,531 & 23/23 & 44/44\\
cli ad2a33837712 & 78,155 & 16,657 & 83,776 & 79,540 & 16,776 & 83,757 & 110/401 & 2/19\\
cli b87e58a1d390 & 26,184 & 6,235 & 68,609 & 32,648 & 6,434 & 68,718 & 14/61 & 3/70\\
cli babd691446b9 & 22,957 & 6,347 & 53,549 & 24,996 & 7,779 & 53,884 & 19/99 & 13/25\\
cli cba9646b8c25 & 86,253 & 32,808 & 103,793 & 90,653 & 36,818 & 102,489 & 134/420 & 30/38\\
\bottomrule
\end{tabular}
\end{table}

\begin{table}[h]
\centering
\caption{Phase 3 per-instance pass rates. Haiku 5.5 uses $k=50$, including the 20 screen samples for W and A1. Sonnet 5.5 uses $k=10$. ``='' marks an arm byte-identical to another arm. ``--='' marks an A1t arm that was not run separately because it is identical to another arm. No rate is given for those entries.}
\label{tab:p3pass}
\resizebox{\textwidth}{!}{%
\begin{tabular}{l ccccccc cccccc}
\toprule
 & \multicolumn{7}{c}{Haiku 5.5} & \multicolumn{6}{c}{Sonnet 5.5}\\
\cmidrule(lr){2-8}\cmidrule(lr){9-14}
instance & W & Wd & A1 & A1r & Wdr & Wpad & A1t & W & Wd & A1 & A1r & Wdr & Wpad\\
\midrule
caddy 30d3bd2a8959 & 0.98 & 0.96 & 0.94 & 1.00 & 1.00 & 1.00 & 0.98= & 1.00 & 1.00 & 1.00 & 1.00 & 1.00 & 1.00\\
caddy 6727b2f3a92b & 0.34 & 0.20 & 0.26 & 0.06 & 0.18 & 0.06 & 0.34= & 0.30 & 0.20 & 0.20 & 0.30 & 0.40 & 0.00\\
caddy c8591a1e46e4 & 0.64 & 0.38 & 0.36 & 0.52 & 0.30 & 0.36 & 0.44 & 0.20 & 0.00 & 0.50 & 0.00 & 0.00 & 0.30\\
cli 2bcdfbb1ac02 & 0.00 & 0.00= & 0.22 & 0.60 & 0.60= & 0.00 & --= & 0.30 & 0.30= & 0.10 & 0.00 & 0.00= & 0.20\\
cli 326faaac8b5c & 0.86 & 0.86= & 0.84 & 0.94 & 0.94= & 0.96 & 0.86= & 1.00 & 1.00= & 1.00 & 1.00 & 1.00= & 1.00\\
cli 3f6a16a9f8c7 & 0.88 & 0.94 & 0.84 & 0.88 & 0.74 & 0.74 & 0.88= & 1.00 & 1.00 & 1.00 & 1.00 & 1.00 & 1.00\\
cli 8f7afbdbdbd6 & 0.98 & 0.98= & 0.92 & 0.62 & 0.62= & 0.28 & 0.46 & 1.00 & 1.00= & 1.00 & 1.00 & 1.00= & 1.00\\
cli ad2a33837712 & 0.46 & 0.60 & 0.58 & 0.46 & 0.74 & 0.40 & 0.46= & 1.00 & 1.00 & 1.00 & 0.70 & 0.60 & 1.00\\
cli b87e58a1d390 & 0.42 & 0.46 & 0.48 & 0.58 & 0.42 & 0.50 & 0.42= & 1.00 & 0.90 & 0.60 & 0.70 & 0.60 & 0.80\\
cli babd691446b9 & 0.34 & 0.60 & 1.00 & 1.00 & 1.00 & 0.16 & --= & 1.00 & 1.00 & 1.00 & 1.00 & 1.00 & 1.00\\
cli cba9646b8c25 & 0.80 & 0.86 & 0.80 & 0.80 & 0.72 & 0.82 & --= & 1.00 & 0.90 & 0.90 & 0.80 & 1.00 & 1.00\\
\bottomrule
\end{tabular}}
\end{table}

\paragraph{Sample sets for Haiku 5.5.}
The Phase 3 amendment does not say whether the Haiku 5.5 contrasts should use all samples or only the post-screen samples. This report takes the post-screen set as primary. In that set, W and A1 use only the samples drawn after the screen ($n=30$ per instance), and the other arms use $n=50$. The choice was not preregistered.

The reason for the choice is that the screen samples were used to select the instances, which creates a risk of a winner's curse. The per-instance changes from screen to post-screen illustrate this. For example, on caddy 6727b2f3a92b W moved from 0.45 to 0.27, and on caddy c8591a1e46e4 A1 moved from 0.45 to 0.30.

Both sample sets are reported in Table~\ref{tab:p3c}.

\begin{table}[h]
\centering
\caption{Phase 3 paired contrasts ($n=11$). ``Ratio'' is the median token ratio, available for the post-screen Haiku table. NI = non-inferior at $\eps = 0.05$ (one-sided 95\% LB $> -0.05$). Sonnet: W and A1 are both at 1.00 on 6/11 instances.}
\label{tab:p3c}
\resizebox{\textwidth}{!}{%
\begin{tabular}{l cccc c ccc c ccc}
\toprule
 & \multicolumn{5}{c}{Haiku 5.5, post-screen (primary)} & \multicolumn{3}{c}{Haiku 5.5, all samples} & & \multicolumn{3}{c}{Sonnet 5.5 (ceiling 6/11)}\\
\cmidrule(lr){2-6}\cmidrule(lr){7-9}\cmidrule(lr){11-13}
contrast & mean [95\% CI] & LB & NI & signs & ratio & mean [95\% CI] & LB & NI & & mean [95\% CI] & LB & signs\\
\midrule
Wd $-$ W & $+0.009$ $[-